% !TEX program = xelatex
% 공통수학2 · Ⅰ. 도형의 방정식 · 03 점과 직선 사이의 거리 · 2/2차시
% 학생용 활동지 (답안 없음)

\documentclass[11pt,a4paper]{article}

\usepackage{kotex}
\usepackage{iftex}
\ifPDFTeX\else
  \setmainhangulfont{Noto Serif CJK KR}
  \setsanshangulfont{Noto Sans CJK KR}
\fi

\usepackage[margin=2.1cm]{geometry}
\usepackage{parskip}
\usepackage{enumitem}
\usepackage{array}
\usepackage{tabularx}
\usepackage{xcolor}
\usepackage{amssymb}
\usepackage{tikz}
\usepackage[most]{tcolorbox}
\usepackage{fancyhdr}
\usepackage{titlesec}

\definecolor{goldc}{HTML}{B07C22}
\definecolor{goldstrong}{HTML}{8A5F16}
\definecolor{indigoc}{HTML}{2B3B57}
\definecolor{indigosoft}{HTML}{6E82A6}
\definecolor{linec}{HTML}{D6DACB}
\definecolor{surface2}{HTML}{E4E8DC}
\definecolor{writeline}{HTML}{C7CCBA}
\definecolor{inksoft}{HTML}{52604F}

\titleformat{\section}{\normalfont\sffamily\bfseries\large\color{indigoc}}{\thesection}{0.6em}{}
\titlespacing{\section}{0pt}{14pt}{6pt}

\newtcolorbox{goalbox}{enhanced, breakable, colback=white, colframe=goldc, boxrule=0.5pt, leftrule=3pt, arc=2pt, left=10pt, right=10pt, top=8pt, bottom=8pt}
\newtcolorbox{sheetbox}{enhanced, breakable, colback=white, colframe=linec, boxrule=0.6pt, arc=3pt, left=11pt, right=11pt, top=9pt, bottom=9pt}

\newcommand{\writespace}[1]{%
  \par\nobreak\vspace{3pt}
  \foreach \n in {1,...,#1} {%
    \noindent\textcolor{writeline}{\rule{\linewidth}{0.4pt}}\par\vspace{0.62cm}%
  }
}
\newcommand{\fillblank}[1]{\makebox[#1]{\hrulefill}}

\pagestyle{fancy}
\fancyhf{}
\renewcommand{\headrulewidth}{0pt}
\fancyfoot[L]{\footnotesize\color{inksoft} 공통수학2 · 2026학년도 2학기 · 지평선고등학교}
\fancyfoot[R]{\footnotesize\color{inksoft} 다음 시간 --- 8차시 중단원 마무리}

\begin{document}

\noindent
\begin{minipage}[b]{0.62\linewidth}
  {\sffamily\footnotesize\color{indigosoft} 공통수학2 · Ⅰ. 도형의 방정식}\\[2pt]
  {\sffamily\bfseries\Large 03 점과 직선 사이의 거리 \textperiodcentered\ 7차시}
\end{minipage}\hfill
\begin{minipage}[b]{0.36\linewidth}
  \raggedleft\small\color{inksoft}
  반\ \fillblank{1.6cm} \quad 번호\ \fillblank{1.6cm}\\[4pt]
  이름\ \fillblank{3.6cm}
\end{minipage}

\vspace{4pt}
\noindent{\color{goldc}\rule{\linewidth}{2.2pt}}
\vspace{10pt}

\begin{goalbox}
\textbf{오늘의 목표}\\
평행한 두 직선 사이의 거리를 구하고, 좌표만으로 삼각형의 넓이를 구하는 방법을 탐구할 수 있다.
\vspace{4pt}

\small\color{inksoft}
$\square$ 자신 있음 \quad $\square$ 조금 헷갈림 \quad $\square$ 처음 봄 \hfill (수업 전 나의 상태에 표시)
\end{goalbox}

\section{개념 정리 --- 평행한 두 직선 사이의 거리}

\begin{sheetbox}
평행한 두 직선 사이의 거리는 한 직선 위의 \fillblank{4cm} 점과 다른 직선 사이의 거리로 구할 수 있다.

\vspace{8pt}
\textbf{연습} --- 평행한 두 직선 $x+2y+1=0$과 $x+2y+6=0$ 사이의 거리를 구해 보자. (힌트: 첫 번째 직선 위의 점 $(1,-1)$을 이용)
\writespace{2}
\end{sheetbox}

\section{탐구 활동 --- 좌표만으로 삼각형의 넓이 구하기}

\begin{sheetbox}
원점 O, 두 점 A$(x_1,y_1)$, B$(x_2,y_2)$를 꼭짓점으로 하는 삼각형 OAB의 넓이 $S$를, 원점과 직선 사이의 거리를 이용하여 구해 보자.

\begin{enumerate}[leftmargin=1.4em, itemsep=4pt]
  \item 변 AB의 길이 $\overline{\text{AB}}$를 $x_1,y_1,x_2,y_2$로 나타내 보자.
  \item 직선 AB의 방정식을 구하고, 원점과 이 직선 사이의 거리를 구해 보자.
  \item 1번과 2번의 결과를 이용하여, 삼각형 OAB의 넓이가
  \[
  S = \dfrac{1}{2}|x_1y_2 - x_2y_1|
  \]
  임을 설명해 보자. (힌트: $S=\dfrac{1}{2}\times\text{밑변}\times\text{높이}$)
\end{enumerate}
\writespace{6}
\end{sheetbox}

\section{연습 문제}

\begin{sheetbox}
\textbf{1.} 직선 $x-3y+2=0$에 평행하고 원점에서의 거리가 $\sqrt{10}$인 직선의 방정식을 구하시오.
\writespace{2}

\vspace{8pt}
\textbf{2.} 다음 직선의 방정식을 구하시오.\\
\hspace*{1.6em}⑴ 직선 $3x+4y-5=0$에 평행하고 원점에서의 거리가 $2$인 직선\\
\hspace*{1.6em}⑵ 직선 $2x-y-4=0$에 수직이고 점 $(2,-1)$에서의 거리가 $4$인 직선
\writespace{4}
\end{sheetbox}

\section{사고의 가시화 --- Connect · Extend · Challenge}

\noindent{\footnotesize\color{inksoft} `점과 직선 사이의 거리' 전체(6$\sim$7차시)를 떠올리며 아래 세 칸을 채워 보자.}\\[6pt]
\begin{tabularx}{\linewidth}{@{}XXX@{}}
\textbf{\footnotesize\color{indigoc} Connect --- 무엇과 연결되나?} &
\textbf{\footnotesize\color{indigoc} Extend --- 어디까지 확장됐나?} &
\textbf{\footnotesize\color{indigoc} Challenge --- 아직 궁금한 것은?} \\[4pt]
\end{tabularx}
\noindent
\begin{minipage}[t]{0.32\linewidth}\writespace{3}\end{minipage}\hfill
\begin{minipage}[t]{0.32\linewidth}\writespace{3}\end{minipage}\hfill
\begin{minipage}[t]{0.32\linewidth}\writespace{3}\end{minipage}

\section{마무리 성찰}

\begin{sheetbox}
\begin{minipage}[t]{0.48\linewidth}
{\footnotesize\color{inksoft} `점과 직선 사이의 거리' 소단원을 마치며 한 문장으로}
\writespace{2}
\end{minipage}\hfill
\begin{minipage}[t]{0.48\linewidth}
{\footnotesize\color{inksoft} 감사 일기 / 유념 한 줄}
\writespace{2}
\end{minipage}
\end{sheetbox}

\end{document}
